\section{Cyclic quantitative contracts from finite weighted certificates}
\label{sec:feedback}

The existing numerical precision-cost construction already handles a finite
acyclic dependency graph: it constructs the rational error budgets and a safe
dyadic schedule internally, with the numerical node count as a variable
\cite{repo-dag}. Replacing its graph by a cycle is not justified by the
contract-cut rule. A cycle must instead be discharged by a finite numerical
invariant. We give such a discharge, useful for coupled implicit equations,
preconditioned polynomial systems, and mutually dependent approximation
blocks. No fixed point, completed real space, or limit is assumed in the finite
theorems below.

\subsection{Data and the certificate}

Fix a numerical dimension $m$. A vector is a finite rational list of length
$m$; a matrix is a finite rational list of $m$ rows of that length. All rational
numbers have positive denominators. Let $A=(a_{ij})$ have nonnegative entries
and let $w=(w_i)$ have positive rational entries. The certificate is the finite
set of inequalities
\begin{equation}\label{eq:half-certificate}
 a_{ij}\geq0,\qquad w_i>0,\qquad
 2\sum_{j<m} a_{ij}w_j\leq w_i\quad(i<m).
\end{equation}
Rational weights can be multiplied by a common positive denominator, so
positive integer weights lose no generality. The checker therefore needs only
integer comparisons, rational arithmetic, and bounded sums. Neither
eigenvalues nor an extremum over a real sphere is involved.

For rational vectors define
\[
 d_w(x,y)=\max_{i<m}\frac{|x_i-y_i|}{w_i},
\]
with value $0$ when $m=0$. This maximum is a bounded rational scan. In the
object theory it is usually more convenient to replace $d_w(x,y)\leq r$ by
the Boolean fold of the inequalities $|x_i-y_i|\leq rw_i$ for $i<m$.

Let $D$ be a rational domain given by a quantifier-free decidable predicate,
and let $T$ be one fixed primitive-recursive rational-vector program, possibly
with numerical data parameters and variable dimension. Its required local
contracts are domain preservation and the componentwise gain estimate
\begin{align}
 x\in D&\Longrightarrow T(x)\in D,\label{eq:T-domain}\\
 |T_i(x)-T_i(y)|&\leq\sum_{j<m}a_{ij}|x_j-y_j|
 \quad(x,y\in D, i<m).\label{eq:T-gain}
\end{align}
The contracts are assumptions to be discharged for the chosen program.
For the polynomial applications later in the report they are finite
coefficient and box estimates; the present theorem does not infer them from
the mere presence of a matrix in a file.

\begin{lemma}[Weighted finite contraction]\label{lem:weighted-contraction}
Under \eqref{eq:half-certificate} and \eqref{eq:T-gain},
\[
 d_w(Tx,Ty)\leq\tfrac12 d_w(x,y)\qquad(x,y\in D).
\]
More directly, if $|x_j-y_j|\leq r w_j$ for every $j<m$ and $r\geq0$,
then $|T_i(x)-T_i(y)|\leq(r/2)w_i$ for every $i<m$.
\end{lemma}
\begin{proof}
Multiplication by $a_{ij}\geq0$ preserves the input inequalities. Their finite
sum gives
\[
 |T_i(x)-T_i(y)|\leq r\sum_{j<m}a_{ij}w_j\leq\tfrac r2 w_i.
\]
This is already the desired componentwise conclusion. For the norm notation,
take $r$ to be the explicitly computed finite maximum. The empty-dimensional
case has no component assertions and is immediate.
\end{proof}

Weights matter even when all entries are nonnegative. For example,
\[
 A=\begin{pmatrix}0&3\\1/16&0\end{pmatrix},\qquad
 w=\begin{pmatrix}8\\1\end{pmatrix}
\]
satisfy the certificate although the first ordinary row sum is $3$.
The feedback is cyclic; an ordinary unweighted contraction estimate would
miss this certificate.

\subsection{Finite iteration constructs a name}

Define the orbit by the primitive recursion
\[
 x_0=c,\qquad x_{n+1}=T(x_n).
\]
Suppose $c\in D$ and a nonnegative rational $B$ satisfies
$|T_i(c)-c_i|\leq Bw_i$ for all $i<m$.

\begin{theorem}[Finite tail estimate and a primitive-recursive name]
\label{thm:feedback-name}
Every orbit element lies in $D$. For all natural $n,j$ and every $i<m$,
\begin{align}
 |x_{n+1,i}-x_{n,i}|&\leq B2^{-n}w_i,\label{eq:feedback-inc}\\
 |x_{n+j,i}-x_{n,i}|&\leq
 2B2^{-n}(1-2^{-j})w_i\leq 2B2^{-n}w_i.
 \label{eq:feedback-tail}
\end{align}
Let $W\geq1$ be an integer upper bound for every $w_i$, and choose a natural
$c_0$ with $2^{c_0}\geq2BW$. Then
$X(k)=x_{k+c_0}$ is a primitive-recursive rational name satisfying
\[
 |X_i(\ell)-X_i(k)|\leq2^{-k}\quad(k\leq\ell, i<m).
\]
Moreover $T(X(k))-X(k)$ tends to zero with the explicit bound
$B2^{-(k+c_0)}w_i$.
\end{theorem}
\begin{proof}
Domain preservation proves the first assertion by induction. The first
increment is bounded by $B$; the contraction lemma halves each succeeding
increment, proving \eqref{eq:feedback-inc} by induction on $n$.
For fixed $n$, telescope the $j$ increments and use their bounds. The finite
geometric identity
\[
 \sum_{h<j}2^{-h}=2(1-2^{-j})
\]
is proved by induction on $j$ and gives \eqref{eq:feedback-tail}. Finally put
$n=k+c_0$ and $j=\ell-k$. Since $w_i\leq W$ and $2BW\leq2^{c_0}$, the
right side is at most $2^{-k}$. The residual is the next increment. Every
statement used here is about finitely many rational values; a real completion
may subsequently interpret the constructed name, but does not produce it.
\end{proof}

There is no unbounded logarithm search in the clock. If $J$ is the least
integer at least $\max(1,2BW)$, a scan through $0,\ldots,J$ finds $c_0$
with $2^{c_0}\geq J$, because $2^J\geq J$ for $J\geq1$.
The resulting least such exponent is primitive recursive. One may instead
reuse the repository's binary-length routine. The uniformity is in the
numerical coefficients, dimension, weights, and requested precision, for one
fixed evaluator template. It is not uniform evaluation of arbitrary
primitive-recursive source codes.

An invariant box can itself be supplied by a finite certificate. If
$D=\{x:|x_i-c_i|\leq Rw_i\ (i<m)\}$, the gain estimate holds on this box,
and $|T_i(c)-c_i|\leq(R/2)w_i$, then
\[
 |T_i(x)-c_i|\leq |T_i(x)-T_i(c)|+|T_i(c)-c_i|
 \leq (R/2)w_i+(R/2)w_i=Rw_i.
\]
Thus domain preservation follows from the center residual and the gain
certificate. A strict spare margin allows subsequent rounding errors.

For a general rational domain $D$, the name need not represent a point of $D$ itself: $D=(0,1)$ and $T(x)=x/2$ give a counterexample. Identifying its completion as a fixed point in the declared domain requires a completion-closed domain and a continuous extension of $T$. The closed boxes and polynomial maps used below have these properties. The finite inequalities above need neither condition.

\subsection{Residual comparison and data perturbations}

The next result is useful even when the orbit is not explicitly used.

\begin{theorem}[A posteriori finite comparison]\label{thm:feedback-residual}
Let $x,y\in D$ and suppose
$d_w(Tx,x)\leq r$ and $d_w(Ty,y)\leq s$, with $r,s\geq0$.
Then
\[
 d_w(x,y)\leq2(r+s).
\]
More generally, let $S$ be a second map, $y\in D$, and suppose
$d_w(Sy,y)\leq s$ and $d_w(Ty,Sy)\leq\delta$.
The conclusion becomes
\begin{equation}\label{eq:perturbed-residual}
 d_w(x,y)\leq2(r+s+\delta).
\end{equation}
Only $T$ needs the contraction certificate for this comparison.
\end{theorem}
\begin{proof}
For $m=0$ all discrepancies vanish. Otherwise the bounded rational maximum
$E=d_w(x,y)$ is attained at a computed index $i_0$. Three triangle
inequalities, with one extra term in the perturbed case, give
\[
 E\leq r+\tfrac12 E+s+\delta.
\]
Here division by $w_{i_0}$ is permitted because it is a positive rational
from the certificate. Subtracting $E/2$ proves
$E\leq2(r+s+\delta)$. No maximum of a represented real function and no
excluded-middle decision about real equality occurs.
\end{proof}

Applied to the names of Theorem~\ref{thm:feedback-name}, this finite estimate
gives uniqueness at the represented level. To spell out the passage, let
two rational names have residual bounds converging to zero with supplied
primitive-recursive moduli. At any requested precision, choose the two
indices where $2(r+s)$ is below that precision; the displayed finite theorem
bounds their distance. Their names are therefore equivalent. With parameter
perturbations the same argument gives a computable continuity modulus for
the solution, provided the same domain and weight certificate remain valid.
This assertion is local to the certified branch. It neither selects a branch
across a singularity nor claims global uniqueness of a polynomial system.

\subsection{Inexact iteration and genuinely simultaneous error budgets}

Suppose two finite trajectories $u_0,\ldots,u_n$ and
$v_0,\ldots,v_n$ lie in $D$ and have local defects
\[
 d_w(u_{j+1},T(u_j))\leq\varepsilon_j,
 \qquad d_w(v_{j+1},T(v_j))\leq\eta_j.
\]
Contraction and the triangle inequality yield
\[
 E_{j+1}\leq\tfrac12E_j+\varepsilon_j+\eta_j,
 \qquad E_j=d_w(u_j,v_j).
\]
Induction gives the finite convolution bound
\begin{equation}\label{eq:feedback-inexact}
 E_n\leq 2^{-n}E_0+
 \sum_{j<n}2^{-(n-1-j)}(\varepsilon_j+\eta_j).
\end{equation}
The base case is equality. For the step, multiply the preceding bound by
$1/2$ and add the new defect; this shifts every exponent and appends precisely
the last summand. This proof handles all coupled coordinates at once.

For an exact trajectory $v$ and a rounded trajectory $u$ with the same initial
point, take $\varepsilon_j\leq\varepsilon4^{-j}$ and $\eta_j=0$.
For $n\geq1$,
\[
 \sum_{j<n}2^{-(n-1-j)}\varepsilon4^{-j}
 =\varepsilon2^{-(n-1)}\sum_{j<n}2^{-j}
 \leq4\varepsilon2^{-n}.
\]
The same upper bound holds at $n=0$. Combining this with
\eqref{eq:feedback-tail} proves
\[
 d_w(u_{n+j},u_n)\leq(2B+8\varepsilon)2^{-n}.
\]
Consequently a precision schedule $p_j=2j+c_1$ for a rounded fixed rational
kernel supplies a name with an explicit clock. More precisely, for $m>0$ and $\varepsilon>0$, if each coordinate rounding error is at most $2^{-p_j}$, choose $c_1$ by bounded search so that $2^{-c_1}\le\varepsilon\min_i w_i$. Then the weighted defect is at most $\varepsilon4^{-j}$. If $\varepsilon=0$, use exact iteration; if $m=0$, the trajectory is empty. The proof requires the rounded
iterates to stay in the domain; this must be checked from a spare invariant-box
margin or enforced by a certified projection. The tail estimate alone does
not imply domain preservation.

\subsection{A finite Neumann producer}

For the nonnegative matrix $A$ in the certificate and any rational right side
$b$, define
\[
 z_0=0,\qquad z_{n+1}=Az_n+b.
\]
These are the actual finite sums
$z_n=\sum_{j<n}A^j b$, as induction verifies. The exact residual identity is
\[
 (I-A)z_n=b-A^n b.
\]
If $|b_i|\leq Bw_i$, the contraction lemma for the linear map gives
$|(A^n b)_i|\leq B2^{-n}w_i$. Hence the recurrence produces a solution name
to $(I-A)z=b$ with an explicit residual modulus and the same finite tail
bound as above. This is an effective producer for a cyclic linear block.
It uses neither a pre-existing inverse nor a convergence theorem from
classical analysis. Arbitrary finite matrix dimensions are handled by the
same bounded rational sum program.

\subsection{Object-theory induction and the Lean boundary}

All row estimates, domain tests, and coordinate comparisons have a decidable
rational meaning. Clear their positive denominators and encode a bounded
universal assertion by a primitive-recursive violation sum. For example the
increment invariant is the equation
\[
 \sum_{i<m}\bigl[\,2^n|x_{n+1,i}-x_{n,i}|>Bw_i\,\bigr]=0,
\]
where the bracket denotes the rational comparison bit after denominator
clearing. This is one quantifier-free equation, even though $m,n$ are numerical
variables. The tail proof uses the corresponding bounded fold over the
coordinate inequalities; the finite geometric identity and convolution
recurrence use ordinary quantifier-free induction. If domain preservation
is needed simultaneously, pair its Boolean invariant with the increment
invariant in the coded state.

The accompanying \nolinkurl{FeedbackReference.lean} gives definitions and
fully written candidate proofs for weighted domination, contraction, finite
increments, finite tails, bounded-maximum witnesses, and residual comparison
over Lean's rational numbers. It has not been compiled in this environment.
It is not an object-theory proof. The separate port plan specifies the
rational/vector codes and the required endpoints of the form
$\Gamma\vdash[\mathrm{PRAAx}]\varphi$. A successful compilation of the
reference file would still leave that object-theory port to be completed.
The exact-rational Python examples check cyclic linear and nonlinear
instances, including a matrix whose unweighted row sum exceeds one; those
checks support the implementation but are not proofs in iPRA.

The advance over acyclic budgets is thus specific: a finite positive-weight
certificate converts a cyclic numerical dependency into a terminating
precision-controlled producer. A cycle of unproved logical contracts remains
unproved; no global reflection or self-supporting proof rule is introduced.
